\documentclass[12pt,a4paper]{article}

% Idioma y codificación (si escribes en español)
\usepackage[utf8]{inputenc}
\usepackage[spanish]{babel}

% Matemáticas y Física
\usepackage{amsmath, amssymb, amsthm}
\usepackage{mathtools}
\usepackage{bm}
\usepackage{physics}
\usepackage{cancel}
\usepackage{bbold}

% Gráficos y Tablas
\usepackage{graphicx}
\usepackage{booktabs}
\usepackage{float}
\usepackage[font=small, labelfont=bf]{caption}
\usepackage{subfigure} 

% Referencias (hyperref siempre debe ser de los últimos)
\usepackage{hyperref}
\usepackage{cleveref}      % colors

% Formato de la página, encabezado y pie de página
\usepackage{fancyhdr}
\setlength{\headheight}{16pt}
\pagestyle{fancy}
\fancyhf{}
\fancyhead[R]{\nouppercase{\leftmark}}
\fancyfoot[R]{\thepage}

\author{Olmo Viviens Blázquez}
\title{Emergencia de la Criticidad en Redes Neuronales Mediante el Aprendizaje}
\date{\today}

% Inicio del documento
\begin{document}

% Portada
\maketitle

\newpage

% Índice o tabla de contenidos
\tableofcontents

\newpage

\section{Introducción}
\begin{enumerate}
  \item Sistemas críticos
  \item Aprendizaje
  \item Criticidad autoorganizada
\end{enumerate}

\section{Fundamento teórico}
\subsection{Modelo anti-hebbiano}
Se va a emplear un modelo lineal de red neuronal en aras de simplificar el proceso de análisis y la interpretación de resultados, a pesar de la diferencia de la representación empleada con la realidad biológica. De este modo la actividad de la neurona $i$ perteneciente a un sistema de $N$ neuronas conectadas entre sí evoluciona según la ecuación \eqref{eq:evolución actividad}
\begin{equation}
  \dot{x}_i = \sum_{j=1}^N w_{ij}x_j
  \label{eq:evolución actividad}
\end{equation}
donde $w_{ij}$ son los elementos de la matriz de conexiones $W$. Esta ecuación puede ser escrita en forma matricial como 

\begin{equation}
  \dot{x} = Wx
\end{equation}

Las conexiones entre neuronas evolucionan en el tiempo según la ecuación \eqref{eq:evolución conexiones}

\begin{equation}
  \dot{w}_{ij} = \alpha(\delta_{ij} - x_i x_j)
  \label{eq:evolución conexiones}
\end{equation}
que en forma matricial se expresa como 

\begin{equation}
  \dot{W} = \alpha\left(\mathbb{1}_N - xx^T\right)
\end{equation}
donde $\alpha$ es el parámetro de aprendizaje $\ll 1$ por lo general.

\subsection{Aprendizaje antisimétrico}

\begin{equation}
  \Delta_l = xy^T - yx^T
\end{equation}

\subsection{Estabilidad al borde del caos}
\section{Resultados}
\subsection{Aprendizaje antihebbiano}
\cite{magnasco2009}

\subsection{Almacenamiento de recuerdos imaginarios}
\cite{susman2019}

\newpage

\bibliographystyle{unsrt}
\bibliography{bibliografia}
\end{document}